Um zu zeigen, dass $2 E(\mathbf{k}) \delta^{(3)}\left(\mathbf{k}-\mathbf{k}^{\prime}\right)$ eine Lorentz-invariante Verteilung ist, beginnen wir mit dem Lorentz-invarianten Ausdruck

\[
\frac{d^4 k}{(2 \pi)^3} \delta\left(k^2-m^2\right) \theta\left(k_0\right).
\]

Dieser Ausdruck ist Lorentz-invariant, da das Maß $d^4 k$ und die Dirac-Delta-Funktion $\delta\left(k^2-m^2\right)$ invariant sind, und die Heaviside-Funktion $\theta\left(k_0\right)$ sicherstellt, dass nur positive Energien berücksichtigt werden.

Wir verwenden die Identität für die Delta-Funktion

\[
\delta\left(k^2-m^2\right) = \frac{1}{2|E(\mathbf{k})|} \left( \delta\left(k_0 - E(\mathbf{k})\right) + \delta\left(k_0 + E(\mathbf{k})\right) \right),
\]

wobei $E(\mathbf{k}) = \sqrt{\mathbf{k}^2 + m^2}$ die positive Energie ist. Da $\theta(k_0)$ nur den positiven Energieanteil berücksichtigt, bleibt nur der Term $\delta\left(k_0 - E(\mathbf{k})\right)$ übrig. Somit reduziert sich der Ausdruck auf

\[
\frac{d^4 k}{(2 \pi)^3} \frac{1}{2 E(\mathbf{k})} \delta\left(k_0 - E(\mathbf{k})\right).
\]

Durch Integration über $k_0$ erhalten wir

\[
\frac{d^3 \mathbf{k}}{(2 \pi)^3} \frac{1}{2 E(\mathbf{k})}.
\]

Um die Lorentz-Invarianz von $2 E(\mathbf{k}) \delta^{(3)}\left(\mathbf{k}-\mathbf{k}^{\prime}\right)$ zu zeigen, multiplizieren wir mit $2 E(\mathbf{k})$, was ergibt

\[
\frac{d^3 \mathbf{k}}{(2 \pi)^3} \delta^{(3)}\left(\mathbf{k}-\mathbf{k}^{\prime}\right).
\]

Da das Maß $d^3 \mathbf{k}$ und die dreidimensionale Delta-Funktion $\delta^{(3)}\left(\mathbf{k}-\mathbf{k}^{\prime}\right)$ unter räumlichen Drehungen und Boosts invariant sind, ist der gesamte Ausdruck Lorentz-invariant. Dies zeigt, dass $2 E(\mathbf{k}) \delta^{(3)}\left(\mathbf{k}-\mathbf{k}^{\prime}\right)$ eine Lorentz-invariante Verteilung ist.